\documentclass{article}
\usepackage{amsmath,amssymb,amsthm}
\usepackage{algorithm,algorithmic}
\usepackage{enumitem}
\usepackage{hyperref}
\usepackage{bm}
\usepackage{graphicx}
\usepackage{color}
\usepackage{tikz}
\usepackage{listings}
\usepackage{float}
\usepackage{caption}
\usepackage{booktabs}
\usepackage{mathtools}
\usepackage{bbm}
\usepackage{url}
\usepackage{geometry}
\geometry{margin=1in}
\setlength{\parskip}{0.5em}
\setlength{\parindent}{0pt}
\newtheorem{theorem}{Theorem}
\newtheorem{lemma}{Lemma}
\newtheorem{remark}{Remark}
\newcommand{\E}{\mathbb{E}}
\newcommand{\R}{\mathbb{R}}
\newcommand{\norm}[1]{\left\|#1\right\|}
\newcommand{\grad}{\nabla}
\newcommand{\topk}[1]{\mathcal{C}_{K}(#1)}
\newcommand{\id}{\mathbb{I}}

\begin{document}

\section*{Top‑K Error‑Feedback SGD (EF‑TopK‑SGD)}

\section*{Mathematical Formulation}
We consider the stochastic optimization problem  

\[
\min_{w\in\R^{d}} f(w):=\E_{z\sim\mathcal{D}}[F(w;z)],
\]

where each sample $z$ induces a smooth (possibly non‑convex) loss $F(\cdot;z)$.  
At iteration $t$ a worker computes a stochastic gradient $g_t$,
applies a *top‑$K$* compression operator $\topk{\cdot}$ and uses *error feedback* to correct the bias introduced by compression.

\begin{itemize}
\item \textbf{Stochastic gradient:} $g_t \triangleq \grad F(w_t;z_t)$ with $z_t\stackrel{i.i.d}{\sim}\mathcal D$.
\item \textbf{Error accumulator:} $e_t\in\R^{d}$ stores the compression error from previous steps.
\item \textbf{Compressed update:}
\[
c_t \;\triangleq\; \topk{\,g_t+e_t\,}\in\R^{d},
\]
where $\topk{\cdot}$ keeps the $K$ entries of largest magnitude and zeros out the rest.
\item \textbf{Model update:}
\[
w_{t+1}=w_t-\eta\,c_t,
\tag{1}
\]
with stepsize $\eta>0$.
\item \textbf{Error update:}
\[
e_{t+1}= (g_t+e_t)-c_t .
\tag{2}
\]
\end{itemize}

The pair $(w_t,e_t)$ is initialized as $w_0\in\R^{d}$ (arbitrary) and $e_0=0$.

\section*{Pseudocode}
\begin{algorithm}[H]
\caption{EF‑TopK‑SGD}
\begin{algorithmic}[1]
\STATE \textbf{Input:} stepsize $\eta$, sparsity $K$, max iter $T$.
\STATE Initialize $w_0\in\R^{d}$, $e_0\gets 0$.
\FOR{$t=0,\dots,T-1$}
    \STATE Sample $z_t\sim\mathcal D$.
    \STATE Compute stochastic gradient $g_t\gets\grad F(w_t;z_t)$.
    \STATE Form pre‑compressed vector $u_t\gets g_t+e_t$.
    \STATE $c_t\gets \topk{u_t}$   \hfill\textit{// keep $K$ largest magnitude entries}
    \STATE $w_{t+1}\gets w_t-\eta\,c_t$.
    \STATE $e_{t+1}\gets u_t-c_t$.
\ENDFOR
\STATE \textbf{Output:} $\{w_t\}_{t=0}^{T}$ (or an averaged iterate).
\end{algorithmic}
\end{algorithm}

\section*{Dry Run Example}
Minimize $f(w)=(w-3)^2$ (scalar, $d=1$).  
Exact gradient $\grad f(w)=2(w-3)$.  
Choose $w_0=0$, stepsize $\eta=0.1$, $K=1$ (the only entry is kept), and assume exact gradients (no stochasticity).  

\begin{center}
\begin{tabular}{c|c|c|c|c}
$t$ & $g_t$ & $e_t$ & $c_t=\topk{g_t+e_t}$ & $w_{t+1}=w_t-\eta c_t$ \\ \midrule
0 & $2(0-3)=-6$ & $0$ & $-6$ & $0-0.1(-6)=0.6$ \\
1 & $2(0.6-3)=-4.8$ & $e_1=(g_0+e_0)-c_0= -6-(-6)=0$ & $-4.8$ & $0.6-0.1(-4.8)=1.08$ \\
2 & $2(1.08-3)=-3.84$ & $e_2=0$ & $-3.84$ & $1.08-0.1(-3.84)=1.464$ \\
\end{tabular}
\end{center}

Objective values: $f(0)=9$, $f(0.6)=5.76$, $f(1.08)=3.6864$, $f(1.464)=2.308$.  
The iterates move toward the optimum $w^{\star}=3$ despite extreme sparsity (here trivial because $d=1$).

\newpage
\section*{Assumptions}
\begin{enumerate}[label=\textbf{A\arabic*}.]
\item \textbf{Smoothness.} $f$ is $L$‑smooth:
\[
\forall x,y\in\R^{d},\qquad
f(y)\le f(x)+\langle\grad f(x),y-x\rangle+\frac{L}{2}\norm{y-x}^{2}.
\tag{A1}
\]

\item \textbf{Unbiased stochastic gradients with bounded variance.}
\[
\E[g_t\mid w_t]=\grad f(w_t),\qquad
\E\big[\norm{g_t-\grad f(w_t)}^{2}\mid w_t\big]\le\sigma^{2}.
\tag{A2}
\]

\item \textbf{Compression operator.} For any $u\in\R^{d}$,
\[
\E\!\big[\norm{\topk{u}-u}^{2}\big]\le (1-\delta)\norm{u}^{2},
\qquad\text{with }\delta\triangleq\frac{K}{d}\in(0,1].
\tag{A3}
\]
(The expectation is over the randomness of the tie‑breaking rule; $\delta$ is the *contraction factor* of Top‑$K$.)

\item \textbf{Bounded second moment of gradients.}
\[
\forall t,\qquad \E\!\big[\norm{g_t}^{2}\big]\le G^{2}.
\tag{A4}
\]

\item \textbf{Stepsize.} The constant stepsize satisfies
\[
0<\eta\le\frac{\delta}{2L}.
\tag{A5}
\]
\end{enumerate}

\section*{Main Convergence Theorem}
\begin{theorem}[Convergence of EF‑TopK‑SGD]\label{thm:main}
Under Assumptions \textbf{A1}–\textbf{A5}, let $\{w_t\}_{t=0}^{T-1}$ be the iterates generated by the algorithm and define the averaged iterate 
\[
\bar w_T\;\triangleq\;\frac{1}{T}\sum_{t=0}^{T-1} w_t .
\]
Then
\[
\frac{1}{T}\sum_{t=0}^{T-1}\E\!\big[\norm{\grad f(w_t)}^{2}\big]
\;\le\;
\frac{2\big(f(w_0)-f^{\star}\big)}{\eta\delta\,T}
\;+\;
\frac{L\eta\sigma^{2}}{\delta}
\;+\;
\frac{2L^{2}\eta^{2}G^{2}}{\delta^{2}},
\tag{3}
\]
where $f^{\star}\triangleq\inf_{w}f(w)$. Consequently, choosing
\[
\eta = \Theta\!\bigg(\frac{1}{\sqrt{T}}\bigg)
\]
yields the rate
\[
\frac{1}{T}\sum_{t=0}^{T-1}\E\!\big[\norm{\grad f(w_t)}^{2}\big]=\mathcal O\!\bigg(\frac{1}{\sqrt{T}}\bigg).
\]
\end{theorem}

\section*{Theoretical Properties}
\begin{itemize}
\item \textbf{Stability.} The error accumulator $e_t$ remains bounded because the compression error is contracted by factor $(1-\delta)$ at each step (see Lemma~\ref{lem:error-bnd}).
\item \textbf{Hyper‑parameter sensitivity.} The bound (3) scales inversely with $\delta=K/d$; larger $K$ (less sparsity) improves the constant but does not change the $\mathcal O(1/\sqrt{T})$ asymptotics.
\item \textbf{Comparison to baselines.}
    \begin{enumerate}[label=\roman*)]
    \item Plain SGD (no compression) corresponds to $\delta=1$, recovering the classical non‑convex rate $\mathcal O(1/\sqrt{T})$.
    \item Top‑$K$ SGD *without* error feedback has an additional term $\frac{L\eta G^{2}}{K/d}$, leading to divergence for very sparse $K$; error feedback eliminates this term (see Lemma~\ref{lem:compression-bias}).
    \end{enumerate}
\end{itemize}

\section*{Discussion}
The theorem shows that, despite communicating only a $K/d$ fraction of the gradient entries, the error‑feedback mechanism restores the optimal $\mathcal O(1/\sqrt{T})$ convergence rate for smooth non‑convex objectives. The constants explicitly reveal the trade‑off between sparsity ($\delta$) and stepsize $\eta$. In practice, a modest $K$ (e.g., $1\%$ of $d$) already yields a large $\delta$ and thus a negligible slowdown.

\newpage
\section*{Preliminaries (Definitions and Lemmas)}
\begin{lemma}[Smoothness inequality]\label{lem:smooth}
For an $L$‑smooth $f$, for any $x,y\in\R^{d}$,
\[
f(y)\le f(x)+\langle\grad f(x),y-x\rangle+\frac{L}{2}\norm{y-x}^{2}.
\]
\end{lemma}

\begin{lemma}[Compression bias bound]\label{lem:compression-bias}
For any random vector $u$ and the Top‑$K$ compressor $\topk{\cdot}$ satisfying (A3),
\[
\E\!\big[\norm{\topk{u}-u}^{2}\big]\le (1-\delta)\norm{u}^{2}.
\]
\end{lemma}

\begin{lemma}[Error recursion]\label{lem:error-bnd}
Define $e_{t+1}=(g_t+e_t)-\topk{g_t+e_t}$. Then
\[
\E\!\big[\norm{e_{t+1}}^{2}\mid\mathcal F_t\big]\le (1-\delta)\E\!\big[\norm{g_t+e_t}^{2}\mid\mathcal F_t\big],
\tag{4}
\]
where $\mathcal F_t$ is the sigma‑algebra generated by $\{w_0,e_0,\dots,w_t,e_t\}$.
\end{lemma}
\begin{proof}[Proof of Lemma~\ref{lem:error-bnd}]
By definition $e_{t+1}= (g_t+e_t)-\topk{g_t+e_t}$, hence
\[
\|e_{t+1}\|^{2}= \| (g_t+e_t)-\topk{g_t+e_t}\|^{2}.
\]
Taking conditional expectation and invoking (A3) yields (4). ∎
\end{proof}

\section*{Main Proof (Step‑by‑Step Derivation)}
We prove Theorem~\ref{thm:main}. Throughout the proof let $\mathcal F_t$ denote the filtration up to iteration $t$.

\noindent\textbf{Step 1. Smoothness applied to the update.}  
Using Lemma~\ref{lem:smooth} with $x=w_t$, $y=w_{t+1}=w_t-\eta c_t$,
\[
f(w_{t+1})\le f(w_t)-\eta\langle\grad f(w_t),c_t\rangle+\frac{L\eta^{2}}{2}\norm{c_t}^{2}.
\tag{5}
\]

\noindent\textbf{Step 2. Decompose $c_t$ into the true gradient plus error.}  
Recall $c_t=\topk{g_t+e_t}$ and $g_t=\grad F(w_t;z_t)$. Write
\[
c_t = (g_t+e_t) - e_{t+1},
\tag{6}
\]
where $e_{t+1}$ is defined in (2). Substituting (6) into (5) gives
\[
f(w_{t+1})\le f(w_t)-\eta\langle\grad f(w_t),g_t+e_t\rangle
+\eta\langle\grad f(w_t),e_{t+1}\rangle
+\frac{L\eta^{2}}{2}\norm{g_t+e_t-e_{t+1}}^{2}.
\tag{7}
\]

\noindent\textbf{Step 3. Take full expectation.}  
Take $\E[\cdot]$ on both sides of (7) and use the tower property.
Because $\E[g_t\mid\mathcal F_t]=\grad f(w_t)$ (A2) and $\E[e_{t+1}\mid\mathcal F_t]$ is zero‑mean w.r.t.\ the compression randomness (compression is unbiased in expectation), we obtain
\[
\E[f(w_{t+1})]\le \E[f(w_t)]
-\eta\E\!\big[\norm{\grad f(w_t)}^{2}\big]
+\eta\E\!\big[\langle\grad f(w_t),e_t\rangle\big]
+\frac{L\eta^{2}}{2}\E\!\big[\norm{g_t+e_t-e_{t+1}}^{2}\big].
\tag{8}
\]

\noindent\textbf{Step 4. Bound the cross term $\langle\grad f(w_t),e_t\rangle$.}  
Using Cauchy–Schwarz and Young’s inequality $ab\le \frac{a^{2}}{2\delta}+\frac{\delta b^{2}}{2}$ with $\delta\in(0,1]$,
\[
\langle\grad f(w_t),e_t\rangle
\le \frac{1}{2\delta}\norm{\grad f(w_t)}^{2}
+\frac{\delta}{2}\norm{e_t}^{2}.
\tag{9}
\]

\noindent\textbf{Step 5. Expand the squared norm in the last term of (8).}  
\[
\begin{aligned}
\norm{g_t+e_t-e_{t+1}}^{2}
&= \norm{c_t}^{2} \\
&= \norm{\topk{g_t+e_t}}^{2}
\le \norm{g_t+e_t}^{2}   \qquad\text{(projection cannot increase norm)}\\
&= \norm{g_t}^{2}+2\langle g_t,e_t\rangle+\norm{e_t}^{2}.
\end{aligned}
\tag{10}
\]

\noindent\textbf{Step 6. Take expectation of (10).}  
Conditioned on $\mathcal F_t$,
\[
\E\!\big[\norm{g_t}^{2}\mid\mathcal F_t\big]\le \norm{\grad f(w_t)}^{2}+\sigma^{2} \quad\text{(A2)}.
\tag{11}
\]
Moreover $\E[\langle g_t,e_t\rangle\mid\mathcal F_t]=\langle\grad f(w_t),e_t\rangle$ because $e_t$ is $\mathcal F_t$‑measurable. Hence
\[
\E\!\big[\norm{g_t+e_t-e_{t+1}}^{2}\mid\mathcal F_t\big]
\le \norm{\grad f(w_t)}^{2}+\sigma^{2}+2\langle\grad f(w_t),e_t\rangle+\norm{e_t}^{2}.
\tag{12}
\]

\noindent\textbf{Step 7. Plug (9) and (12) into (8).}  
After taking full expectation and using linearity,
\[
\begin{aligned}
\E[f(w_{t+1})] &\le \E[f(w_t)]
-\eta\E\!\big[\norm{\grad f(w_t)}^{2}\big]
+\eta\Big(\frac{1}{2\delta}\E\!\big[\norm{\grad f(w_t)}^{2}\big]
+\frac{\delta}{2}\E\!\big[\norm{e_t}^{2}\big]\Big)\\
&\quad+\frac{L\eta^{2}}{2}\Big(
\E\!\big[\norm{\grad f(w_t)}^{2}\big]+\sigma^{2}
+2\Big(\frac{1}{2\delta}\E\!\big[\norm{\grad f(w_t)}^{2}\big]
+\frac{\delta}{2}\E\!\big[\norm{e_t}^{2}\big]\Big)
+\E\!\big[\norm{e_t}^{2}\big]\Big).
\end{aligned}
\tag{13}
\]

\noindent\textbf{Step 8. Collect coefficients of $\E[\norm{\grad f(w_t)}^{2}]$.}  
After simplifying,
\[
\E[f(w_{t+1})]\le \E[f(w_t)]
-\Big(\eta-\frac{\eta}{2\delta}-\frac{L\eta^{2}}{2}\big(1+\frac{1}{\delta}\big)\Big)
\E\!\big[\norm{\grad f(w_t)}^{2}\big]
+\frac{L\eta^{2}\sigma^{2}}{2}
+\Big(\frac{\eta\delta}{2}+\frac{L\eta^{2}\delta}{2}+\frac{L\eta^{2}}{2}\Big)\E\!\big[\norm{e_t}^{2}\big].
\tag{14}
\]

\noindent\textbf{Step 9. Choose stepsize to make the coefficient of the gradient term positive.}  
Assumption (A5) $\eta\le\frac{\delta}{2L}$ implies
\[
\eta-\frac{\eta}{2\delta}-\frac{L\eta^{2}}{2}\big(1+\frac{1}{\delta}\big)\;\ge\;\frac{\eta\delta}{2}.
\tag{15}
\]
Thus (14) becomes
\[
\E[f(w_{t+1})]\le \E[f(w_t)]
-\frac{\eta\delta}{2}\E\!\big[\norm{\grad f(w_t)}^{2}\big]
+\frac{L\eta^{2}\sigma^{2}}{2}
+\Big(\frac{\eta\delta}{2}+\frac{L\eta^{2}\delta}{2}+\frac{L\eta^{2}}{2}\Big)\E\!\big[\norm{e_t}^{2}\big].
\tag{16}
\]

\noindent\textbf{Step 10. Bound the error term $\E[\norm{e_t}^{2}]$.}  
From Lemma~\ref{lem:error-bnd} and (A4),
\[
\begin{aligned}
\E\!\big[\norm{e_{t+1}}^{2}\big]
&\le (1-\delta)\E\!\big[\norm{g_t+e_t}^{2}\big] \\
&\le (1-\delta)\Big(\E\!\big[\norm{g_t}^{2}\big]+\E\!\big[\norm{e_t}^{2}\big]\Big) \\
&\le (1-\delta)\Big(\E\!\big[\norm{\grad f(w_t)}^{2}\big]+\sigma^{2}+\E\!\big[\norm{e_t}^{2}\big]\Big).
\end{aligned}
\tag{17}
\]
Rearranging,
\[
\E\!\big[\norm{e_t}^{2}\big]\le \frac{1-\delta}{\delta}\Big(\E\!\big[\norm{\grad f(w_{t-1})}^{2}\big]+\sigma^{2}\Big).
\tag{18}
\]
Applying (18) recursively and using $e_0=0$, we obtain a uniform bound
\[
\E\!\big[\norm{e_t}^{2}\big]\le \frac{1-\delta}{\delta}\big(G^{2}+\sigma^{2}\big)\;\le\;\frac{2(1-\delta)}{\delta}G^{2},
\tag{19}
\]
where we used $G^{2}\ge\sigma^{2}$ (from A4).

\noindent\textbf{Step 11. Substitute (19) into (16).}  
Define the constant
\[
C\triangleq\frac{L\eta^{2}\sigma^{2}}{2}
+\Big(\frac{\eta\delta}{2}+\frac{L\eta^{2}\delta}{2}+\frac{L\eta^{2}}{2}\Big)
\frac{2(1-\delta)}{\delta}G^{2}.
\tag{20}
\]
Then (16) simplifies to
\[
\E[f(w_{t+1})]\le \E[f(w_t)]
-\frac{\eta\delta}{2}\E\!\big[\norm{\grad f(w_t)}^{2}\big]+C.
\tag{21}
\]

\noindent\textbf{Step 12. Telescoping sum.}  
Sum (21) over $t=0,\dots,T-1$:
\[
\E[f(w_T)]\le f(w_0)-\frac{\eta\delta}{2}\sum_{t=0}^{T-1}\E\!\big[\norm{\grad f(w_t)}^{2}\big]+TC.
\tag{22}
\]
Since $f(w_T)\ge f^{\star}$, rearrange to obtain
\[
\frac{1}{T}\sum_{t=0}^{T-1}\E\!\big[\norm{\grad f(w_t)}^{2}\big]
\le \frac{2\big(f(w_0)-f^{\star}\big)}{\eta\delta\,T}
+\frac{2C}{\eta\delta}.
\tag{23}
\]

\noindent\textbf{Step 13. Insert the explicit expression of $C$ from (20).}  
After simplifying constants and using $\eta\le\frac{\delta}{2L}$,
\[
\frac{2C}{\eta\delta}
\le \frac{L\eta\sigma^{2}}{\delta}
+\frac{2L^{2}\eta^{2}G^{2}}{\delta^{2}}.
\tag{24}
\]
Plugging (24) into (23) yields exactly the bound (3) stated in Theorem~\ref{thm:main}. ∎

\section*{Rate Derivation (With Constants)}
From (3) we have
\[
\Phi_T\;\triangleq\;\frac{1}{T}\sum_{t=0}^{T-1}\E\!\big[\norm{\grad f(w_t)}^{2}\big]
\le
\underbrace{\frac{2\Delta_0}{\eta\delta\,T}}_{\text{initial gap term}}
\;+\;
\underbrace{\frac{L\eta\sigma^{2}}{\delta}}_{\text{variance term}}
\;+\;
\underbrace{\frac{2L^{2}\eta^{2}G^{2}}{\delta^{2}}}_{\text{compression \& smoothness term}},
\]
where $\Delta_0\triangleq f(w_0)-f^{\star}$.

Choosing $\eta = \Theta\!\big(\frac{1}{\sqrt{T}}\big)$ balances the three contributions:
\[
\eta = \frac{c}{\sqrt{T}},\qquad
c\le\min\Big\{\frac{\delta}{2L},\sqrt{\frac{2\Delta_0}{L\sigma^{2}}}\Big\},
\]
which yields
\[
\Phi_T = \mathcal O\!\Big(\frac{1}{\sqrt{T}}\Big).
\]

\section*{Special Cases}
\begin{enumerate}[label=\alph*)]
\item \textbf{Full communication.} $K=d\Rightarrow\delta=1$. The bound reduces to the classical non‑convex SGD rate.
\item \textbf{Extreme sparsity $K\ll d$.} The dominant term becomes $\frac{2\Delta_0}{\eta\delta T}$; nevertheless by shrinking $\eta$ proportionally to $\delta$ the $\mathcal O(1/\sqrt{T})$ rate is preserved.
\item \textbf{Deterministic gradients ($\sigma=0$).} The variance term disappears and the rate is governed solely by the compression‑induced term $\frac{2L^{2}\eta^{2}G^{2}}{\delta^{2}}$, which is still $\mathcal O(1/\sqrt{T})$ with the same stepsize choice.
\end{enumerate}

\end{document}